\documentclass[../../main.tex]{subfiles}
\begin{document}

{\bf [解]}

(1)

% 視点の設定 (tdplot_set_coordinates{r}{theta}{phi})
\tdplotsetmaincoords{70}{60}

\begin{figure}[htb]
\begin{center}
 \begin{tikzpicture}[tdplot_main_coords, scale=1.8]

    % --- 1. 各頂点・中点の座標定義 (計算値に基づく) ---
    % sqrt(3) ≒ 1.732
    \pgfmathsetmacro{\sqrttwo}{1.4142}
    \pgfmathsetmacro{\sqrtthree}{1.7320}

    % 共通両垂線の足 (中点 E, F)
    \coordinate (E) at (-0.5, -0.5, 0);
    \coordinate (F) at (0.5, 0.5, 1);

    % 直線 l1 上の頂点 A, B
    \coordinate (A) at ({-0.5 + \sqrtthree/2}, {-0.5 - \sqrtthree/2}, 0);
    \coordinate (B) at ({-0.5 - \sqrtthree/2}, {-0.5 + \sqrtthree/2}, 0);

    % 直線 l2 上の頂点 C, D
    \coordinate (C) at ({0.5 + \sqrtthree/2}, {0.5 + \sqrtthree/2}, {1 - \sqrtthree});
    \coordinate (D) at ({0.5 - \sqrtthree/2}, {0.5 - \sqrtthree/2}, {1 + \sqrtthree});

    % --- 2. 座標軸の描画 ---
    \draw[->, thick, gray] (-2,0,0) -- (2.5,0,0) node[anchor=north east]{$x$};
    \draw[->, thick, gray] (0,-2,0) -- (0,2.5,0) node[anchor=north west]{$y$};
    \draw[->, thick, gray] (0,0,-1) -- (0,0,3) node[anchor=south]{$z$};
    \node[anchor=north east, gray] at (0,0,0) {$O$};

    % --- 3. 直線 l1, l2 の描画 ---
    % 直線 l1: A, B を通る
    \draw[dashed, blue!60, line width=0.8pt] 
        ($ (A) ! -0.4 ! (B) $) -- ($ (B) ! -0.4 ! (A) $) 
        node[anchor=south west] {$\ell_1$};

    % 直線 l2: C, D を通る
    \draw[dashed, red!60, line width=0.8pt] 
        ($ (C) ! -0.3 ! (D) $) -- ($ (D) ! -0.3 ! (C) $) 
        node[anchor=north east] {$\ell_2$};

    % --- 4. 正四面体 ABCD の面と辺の描画 ---
    % 奥側の面・隠線（点線）
    \draw[thick, dashed, opacity=0.7] (A) -- (B);

    % 手前側の辺（実線）
    \draw[thick, line join=round] (A) -- (C) -- (B) -- (D) -- cycle;
    \draw[thick] (C) -- (D);

    % --- 5. 共通両垂線 EF (直交関係) ---
    \draw[ultra thick, green!50!black] (E) -- (F);
    % \node[anchor=south, green!40!black] at ($ (E)!0.5!(F) $) {$EF$};

    % --- 6. 垂直記号の補助線描画 ---
    % Eでの垂直 (l1 ⊥ EF)
    \draw[thin, green!50!black] ($ (E)!0.15cm!(A) $) -- ($ (E)!0.15cm!(A) + (E)!0.15cm!(F) - (E) $) -- ($ (E)!0.15cm!(F) $);
    % Fでの垂直 (l2 ⊥ EF)
    \draw[thin, green!50!black] ($ (F)!0.15cm!(D) $) -- ($ (F)!0.15cm!(D) + (F)!0.15cm!(E) - (F) $) -- ($ (F)!0.15cm!(E) $);

    % --- 7. 各点のプロットとラベル指定 ---
    \fill[black] (A) circle (1.2pt) node[anchor=north west] {$A$};
    \fill[black] (B) circle (1.2pt) node[anchor=south east] {$B$};
    \fill[black] (C) circle (1.2pt) node[anchor=north west] {$C$};
    \fill[black] (D) circle (1.2pt) node[anchor=south west] {$D$};

    \fill[green!50!black] (E) circle (1.2pt) node[anchor=south west] {$E$};
    \fill[green!50!black] (F) circle (1.2pt) node[anchor=east] {$F$};

\end{tikzpicture}
\caption{題意の四面体ABCDの様子}
\end{center}
\end{figure}

E,Fも直線$\ell_1,\ell_2$上にあるので$E(t,-1-t,0)$，$F(u,u,-2u+2)$（$t,u\in\mathbb{R}$）とおく．
正四面体の性質から線分$EF$は$\ell_1,\ell_2$と垂直である．$\ell_1,\ell_2$の方向ベクトルを$\vec{v}_1, \vec{v}_2$とすると
\begin{align}
 &\vec{v}_1 = \begin{pmatrix}-1\\1\\0\end{pmatrix} &\vec{v}_2= \begin{pmatrix}1\\1\\-2\end{pmatrix}
\end{align}
であり，
\begin{align}
\overrightarrow{EF}=\begin{pmatrix}u-t\\u+t+1\\-2u+2\end{pmatrix}
\end{align}
だから
\begin{align}
 \vec{v}_1\cdot \vec{EF} = 0 \\
 \vec{v}_2\cdot \vec{EF} = 0 
\end{align}
に代入して
\begin{align}
&\begin{cases}\displaystyle
(u-t)-(u+t+1)=0 \\
(u-t)+(u+t+1)+4(u-1)=0
\end{cases} \\
\therefore
&\begin{cases}\displaystyle
t=-\dfrac{1}{2} \\
u=\dfrac{1}{2}
\end{cases}
\end{align}
となり，
\begin{align}
E\left(-\dfrac{1}{2},-\dfrac{1}{2},0\right), \qquad F\left(\dfrac{1}{2},\dfrac{1}{2},1\right) \label{eq:1}
\end{align}
が求める座標である．


(2) 
1辺の長さを$d$とする．下図のように正四面体の面（正三角形）BCDを考える．

\begin{figure}[htb]
\begin{center}
\begin{tikzpicture}[scale=1.2]
    % 座標設定 (1辺 d の正三角形)
    % C を原点付近に置き、D を右、B を上へ
    \coordinate (C) at (0, 0);
    \coordinate (D) at (3, 0);
    \coordinate (B) at (1.5, {3*sqrt(3)/2});
    \coordinate (F) at (1.5, 0); % CD の中点

    % 面 BCD (正三角形) の描画
    \draw[thick] (B) -- (C) -- (D) -- cycle;
    
    % 中線 BF
    \draw[thick, blue!80!black] (B) -- (F);

    % 垂直記号 (F における CD ⊥ BF)
    \draw[thin] (F) ++(0.2,0) -- ++(0,0.2) -- ++(-0.2,0);

    % 辺の長さの表示
    \node[left] at ($(B)!0.5!(C)$) {$d$};
    \node[right] at ($(B)!0.5!(D)$) {$d$};
    \node[below] at ($(C)!0.5!(F)$) {$\frac{d}{2}$};
    \node[below] at ($(F)!0.5!(D)$) {$\frac{d}{2}$};
    % \node[right, blue!80!black] at ($(B)!0.5!(F)$) {$|BF| = \frac{\sqrt{3}}{2}d$};

    % 点の描画とラベル
    \fill[black] (B) circle (1.5pt) node[above] {$B$};
    \fill[black] (C) circle (1.5pt) node[below left] {$C$};
    \fill[black] (D) circle (1.5pt) node[below right] {$D$};
    \fill[blue!80!black] (F) circle (1.5pt) node[below] {$F$};
\end{tikzpicture}
\caption{$\triangle BCD$の様子}
\end{center} 
\end{figure}

図より
\begin{align}
 |BF| = \dfrac{\sqrt{3}}{2}d
\end{align}
である．

\begin{figure}[htb]
\begin{center}
\begin{tikzpicture}[scale=1.3]
    % --- 座標定義 ---
    % AB の長さ d = 3 とし、中点を E(0,0) とする
    \coordinate (E) at (0, 0);
    \coordinate (A) at (-1.5, 0);
    \coordinate (B) at (1.5, 0);
    % 高さ |EF| = sqrt(2)/2 * d に相当する比率で配置
    \coordinate (F) at (0, 2.12);

    % --- 三角形 ABF と中線 EF の描画 ---
    \draw[thick] (A) -- (B) -- (F) -- cycle;
    \draw[thick, green!50!black] (E) -- (F); % 共通両垂線 EF

    % 垂直記号 (E における AB ⊥ EF)
    \draw[thin] (E) ++(0.18,0) -- ++(0,0.18) -- ++(-0.18,0);

    % --- 各辺・線分のラベル ---
    \node[below] at ($(A)!0.5!(E)$) {$\frac{d}{2}$};
    \node[below] at ($(E)!0.5!(B)$) {$\frac{d}{2}$};
    % \node[left, green!50!black] at ($(E)!0.5!(F)$) {$|EF|$};
    \node[above left] at ($(A)!0.5!(F)$) {$|AF| = \frac{\sqrt{3}}{2}d$};
    \node[above right, blue!80!black] at ($(B)!0.5!(F)$) {$|BF| = \frac{\sqrt{3}}{2}d$};

    % --- 点のプロットとラベル ---
    \fill[black] (A) circle (1.5pt) node[below left] {$A$};
    \fill[black] (B) circle (1.5pt) node[below right] {$B$};
    \fill[green!50!black] (E) circle (1.5pt) node[below] {$E$};
    \fill[blue!80!black] (F) circle (1.5pt) node[above] {$F$};
\end{tikzpicture}
\caption{$\triangle ABF$の様子}
\end{center}
\end{figure}

次に上図で$\triangle FEB$にピタゴラスの定理を用いて，
\begin{align}
\begin{split}
|BF|^2 &= |EF|^2+|BE|^2 \\
\left(\dfrac{\sqrt{3}}{2}d\right)^2&=|EF|^2+\left(\dfrac{d}{2}\right)^2 \\
\therefore
|EF|&=\dfrac{\sqrt{2}}{2}d \quad(\because d,|EF|>0) \label{eq:2}
\end{split}
\end{align}
を得る．一方で\cref{eq:1}より$|EF|=\sqrt{3}$だから，\cref{eq:2}に代入して
\begin{align}
  d=\sqrt{6} \label{eq:3}
\end{align}
である．


(3) 

点$A$は$\ell_1$上にあるので$A(p,-1-p,0)$とおける．$|AF|$をふた通りで表すことで$p$を求めよう．
まず\cref{eq:1}より|AF|の長さを計算すると
\begin{align}
 |AF|
  &=\sqrt{\left(p-\dfrac{1}{2}\right)^2+\left(-p-\dfrac{3}{2}\right)^2+1} \\
  &=\sqrt{2p^2+2p+\dfrac{7}{2}} \label{eq:4}
\end{align}
である．

次に(2)の図および\cref{eq:3}より
\begin{align}
 |AF|=\dfrac{\sqrt{3}}{2}d=\dfrac{3}{2}\sqrt{2} \label{eq:5}
\end{align}
である．

以上\cref{eq:4,eq:5}より
\begin{align}
 \sqrt{2p^2+2p+\dfrac{7}{2}} &= \dfrac{3}{2}\sqrt{2} \\
 2p^2+2p+\dfrac{7}{2}        &= \dfrac{9}{2} \\
 2p^2+2p-1 &= 0 \\
\therefore
 p &= \dfrac{-1\pm\sqrt{3}}{2}
\end{align}
となる．

これら複合二つが$A,B$に対応するが，題意より$A$の方が$x$座標が大きいので
複号正に対応するのが$A$であって，求める$A$の座標は
\begin{align}
A\left(\dfrac{-1+\sqrt{3}}{2},\dfrac{-1-\sqrt{3}}{2},0\right)
\end{align}
である．

\newpage

\end{document}
